% ================================================================
% ==== FILE: content/m_spaces/m4.tex
% ================================================================

% ==============================
%  Ontology of Continua — Core
%  Meta-Space: M4
%  FULL MODULE — FINAL VERSION
% ==============================

\subsubsection{\texorpdfstring{$M_4$}{M_4} Übersicht}
\label{sec:m4-overview}

$M_4$ ist der Metaraum, der die Zulässigkeitsbedingungen dafür bereitstellt
biologische Kontinua der Ebene $K_4$.
Es ist der erste Metaraum, in dem:

\begin{itemize}
    \item eine stabile \emph{semipermeable Grenze} ist zulässig,
    \item die topologische Unterscheidung \textbf{innen/außen} wird zu einer Achse,
    \item nontrivial \emph{Konzentrations-, elektrische, pH- und Redoxgradienten}
    kann aufrechterhalten werden,
    \item aktive Transport- und Energieumwandlungsprozesse sind erlaubt,
    \item Reaktionsnetzwerke sind begrenzt und räumlich koordiniert,
    \item der Übergang von der chemischen zur protobiologischen Struktur ($K_3 \to K_4$)
          ist zulässig.
\end{itemize}

$M_4$ führt die geometrischen und energetischen Bedingungen für ein
Kompartimentierung – die bestimmende Eigenschaft früher Zellen.

% ---------------------------------------------------------------
\subsubsection*{1. Zulässiger Konfigurationsraum \texorpdfstring{$\Omega(M_4)$}{\Omega(M_4)}}

Die zulässigen Konfigurationen in $M_4$ erweitern die von $M_3$ um:

\[
\Omega(M_4)
=
\Big\{
(\phi,\ G_{\mathrm{chem}},\ \rho_{\mathrm{in/out}},\
\Delta\mu,\ g_{ij},\ \partial\Omega,\ \Pi)
\ \Big|\ \text{Zulässigkeitsbedingungen gelten}
\Big\}.
\]

Wo:

\begin{itemize}
    \item $\partial\Omega$ – eine semipermeable Membran, die eine kompakte Grenze bildet,
    \item $\rho_{\mathrm{in/out}}$ — Konzentrationsfelder innen und außen,
    \item $\Delta\mu$ – elektrochemische Potentialgradienten über $\partial\Omega$,
    \item $\Pi$ – Permeabilitätstensor (richtungs- und artabhängig),
    \item $G_{\mathrm{chem}}$ – interne RAF-Netzwerke, die funktionsfähig bleiben
          unter Haft.
\end{itemize}

Zulassungsbedingungen:

\begin{enumerate}
    \item \textbf{Membranstabilität}:
    $\partial\Omega$ muss unter Spannung die mechanische Integrität bewahren
    $T_{\mathrm{mem}} < \Theta_{\mathrm{mem}}$.

    \item \textbf{Permeabilitätsbeschränkungen}:
    $\Pi$ muss endlich, anisotrop und artspezifisch sein.

    \item \textbf{Gradientenkompatibilität}:
    Elektrochemische Gradienten müssen Folgendes erfüllen:
    \[
    |\Delta\mu| < \Theta_{\mathrm{grad-max}}.
    \]

    \item \textbf{Volumenendlichkeit}:
    Der interne Bereich muss ein begrenztes Volumen und eine begrenzte Masse haben.

    \item \textbf{RAF-Kompatibilität}:
    Interne Reaktionsnetzwerke müssen topologisch geschlossen bleiben.
\end{enumerate}

$M_4$ bietet somit die minimale Umgebung, in der Protozellen existieren können.

% ---------------------------------------------------------------
\subsubsection*{2. Zulässige Achsen \texorpdfstring{$A(M_4)$}{A(M_4)}}

Es erscheint eine neue Grundachse:

\[
A_{\mathrm{in/out}} = \{ \text{innen},\ \text{außen} \}.
\]

Es erfüllt:

\[
A_{\mathrm{in/out}} \not\subseteq \operatorname{span}(A_1, A_2, A_3)
\quad\Rightarrow\quad
\text{neue Dimension erlaubt}.
\]

Zu den weiteren zulässigen Achsen gehören:

\begin{itemize}
    \item $A_{\mathrm{perm}}$ – Permeabilitätszustände,
    \item $A_{\mathrm{grad}}$ — Gradientenkonfigurationen (pH, Ionen, Ladung),
    \item $A_{\mathrm{vol}}$ — Volumen-/Druckachse,
    \item $A_{\mathrm{energy}}$ – innere Energiezustände
          (ATP-ähnliche Vorläufer, Redox-Pools).
\end{itemize}

Diese Achsen definieren die interne Funktionsgeometrie früher Zellen.

% ---------------------------------------------------------------
\subsubsection*{3. Potentiale \texorpdfstring{$P(M_4)$}{P(M_4)}}

$M_4$ lässt neue Potenzialklassen zu:

\[
P(M_4) =
\big(
U_{\mathrm{mem}},\
U_{\mathrm{grad}},\
U_{\mathrm{osm}},\
U_{\mathrm{Ladung}},\
U_{\mathrm{Redox}},\
U_{\mathrm{aktiv}}
\Big).
\]

Interpretation:

\begin{itemize}
    \item $U_{\mathrm{mem}}$ – Membranbiegung und Oberflächenspannungspotential,
    \item $U_{\mathrm{grad}}$ – in chemischen/elektrischen Gradienten gespeicherte Energie,
    \item $U_{\mathrm{osm}}$ — osmotisches Druckpotential,
    \item $U_{\mathrm{Ladung}}$ — Coulombische und dielektrische Wechselwirkungen,
    \item $U_{\mathrm{redox}}$ — Redox-Energielandschaft,
    \item $U_{\mathrm{active}}$ – Energieumwandlung, die aktiven Transport ermöglicht.
\end{itemize}

Alle Potenziale müssen sein:

\begin{enumerate}
    \item fast überall differenzierbar,
    \item von unten begrenzt,
    \item kompatibel mit Membranstabilitätsbedingungen.
\end{enumerate}

% ---------------------------------------------------------------
\subsubsection*{4. Schwellenwerte \texorpdfstring{$\Theta(M_4)$}{\Theta(M_4)}}

Wichtige biologische Schwellenwerte eingeführt:

\begin{itemize}
    \item $\Theta_{\mathrm{mem}}$ — Membranbruchschwelle,
    \item $\Theta_{\mathrm{grad}}$ – minimaler Gradient, der zur Aufrechterhaltung der Struktur erforderlich ist,
    \item $\Theta_{\mathrm{grad-max}}$ – maximal tolerierbarer Gradient,
    \item $\Theta_{\mathrm{osm}}$ – osmotische Drucktoleranz,
    \item $\Theta_{\mathrm{pH}}$ — Lebensfähigkeitsintervall für den internen pH-Wert,
    \item $\Theta_{\mathrm{redox}}$ – Redox-Kompatibilitätsschwelle.
\end{itemize}

Kritische Bedingung dafür, dass $K_4$ am Leben bleibt:

\[
T_4 <
\min\{
\Theta_{\mathrm{mem}},\
\Theta_{\mathrm{osm}},\
\Theta_{\mathrm{grad-max}},\
\Theta_{\mathrm{pH}},\
\Theta_{\mathrm{Redox}}
\}.
\]

% ---------------------------------------------------------------
\subsubsection*{5. Flüsse \texorpdfstring{$J(M_4)$}{J(M_4)}}

Zulässige Flüsse umfassen:

\[
J(M_4)
=
\big(
J_{\mathrm{diff}},\
J_{\mathrm{perm}},\
J_{\mathrm{Pumpe}},\
J_{\mathrm{Redox}},\
J_{\mathrm{osm}},\
\nabla_i T_4
\Big).
\]

Wo:

\begin{itemize}
    \item $J_{\mathrm{diff}}$ — passive Diffusion innen/außen,
    \item $J_{\mathrm{perm}}$ — Membranpermeationsflüsse,
    \item $J_{\mathrm{pump}}$ — aktiver Transport gegen Gradienten,
    \item $J_{\mathrm{redox}}$ – redoxzyklische Strömungen,
    \item $J_{\mathrm{osm}}$ — osmotische Gleichgewichtsflüsse,
    \item $\nabla_i T_4$ – Strömungen, die durch strukturelle Spannungsgradienten angetrieben werden.
\end{itemize}

Alle Strömungen müssen der Massenerhaltung innerhalb der Membrangrenze genügen.

% ---------------------------------------------------------------
\subsubsection*{6. Zyklen in \texorpdfstring{$M_4$}{M_4}}

$M_4$ ist der Metaraum, in dem biologische Zyklen zulässig werden:

\begin{itemize}
    \item $C_{\mathrm{RAF-core}}$ — reaktionsautokatalytische Zyklen,
    \item $C_{\mathrm{metabolic}}$ – grundlegender Stoffwechselumsatz,
    \item $C_{\mathrm{pump}}$ – pumpengetriebene aktive Transportzyklen,
    \item $C_{\mathrm{redox}}$ — Elektronentransferzyklen,
    \item $C_{\mathrm{buffer}}$ – pH-Pufferzyklen.
\end{itemize}

Voraussetzung für Zyklenstabilität:

\[
\oint_C dA_i \ungefähr 0
\quad\text{und}\quad
T_4 \ \text{begrenzt}.
\]

Diese Zyklen erzeugen die zeitliche Struktur $\tau(K_4)$.

% ---------------------------------------------------------------
\subsubsection*{7. Grenze \texorpdfstring{$\partial\Omega(M_4)$}{\partial\Omega(M_4)}}

Eine Konfiguration nähert sich $\partial\Omega(M_4)$, wenn:

\begin{itemize}
    \item Membranspannung nähert sich der Bruchgrenze:
    \[
    T_{\mathrm{mem}} \to \Theta_{\mathrm{mem}},
    \]
    \item interner pH verlässt das Lebensfähigkeitsintervall,
    \item-Gradienten überschreiten $\Theta_{\mathrm{grad-max}}$,
    \item osmotisches Ungleichgewicht wird instabil,
    \item Einschluss bricht den RAF-Verschluss im Inneren des Abteils.
\end{itemize}

Das Überqueren von $\partial\Omega(M_4)$ zerstört das protobiologische Kontinuum $K_4$.

% ---------------------------------------------------------------
\subsubsection*{8. Beziehung zu \texorpdfstring{$M_3$}{M_3} und \texorpdfstring{$M_5$}{M_5}}

Der Übergang $K_3 \to K_4$ ist zulässig, wenn:

\[
A_{\mathrm{in/out}} \in A(M_4),
\qquad
\Omega(K_4) \subseteq \Omega(M_4),
\qquad
T_3 \ge\Theta_{\mathrm{dim}}.
\]

Der Übergang zu $K_5$ erfordert:

\[
A_{\mathrm{exc}} \in A(M_5),
\qquad
\Updelta V\ge\Theta_{\mathrm{exc}},
\qquad
\text{Kanäle und elektrische Leitung zulässig}.
\]

Somit ist $M_4$ der Metaraum von:

\begin{itemize}
    \item Membranstruktur,
    \item-Unterteilung,
    \item Gradienten und osmotische Physik,
    \item Energieumwandlung,
    \item früher Stoffwechsel,
    \item minimale biologische Stabilität.
\end{itemize}

Es ist die strukturelle Grundlage für protozelluläres Leben.
